\begin{textsample}{POS Dim 4 – 2021\_c5 – Score 1.00 – 2021\_c5/enade\_2021\_des\_64eae85373.txt}  \label{ex:f4_pos_217}
CAPÍTULO 5 DISTRIBUIÇÃO DOS CONCEITOS 5.1 PANORAMA NACIONAL DA DISTRIBUIÇÃO DOS CONCEITOS O Conceito Enade é um indicador de qualidade que avalia os cursos por intermédio dos desempenhos dos estudantes a partir dos resultados do Enade.
Os valores possíveis vão de 1 ( pior situação ) a 5 ( melhor situação ).
Os cursos com apenas um ou sem nenhum concluinte participante não obtêm o Conceito Enade, ficando Sem Conceito ( SC ).
Na Tabela 5.1, são apresentadas a quantidade e a distribuição de cursos de Design participantes do Enade/2021, por faixa de conceito e Grande Região.
Enfatiza -se, mais uma vez, que a diferença entre os cursos tabulados neste capítulo e no Capítulo 2 corresponde aos cursos que não foram avaliados, ou seja, àqueles, em princípio, sem alunos concluintes que tenham sido inscritos para a prova.
Observando -se os dados da Tabela 5.1, nota -se que, dos 109 cursos participantes, 36 ( 33,0 \% ) classificaram -se com conceito 4, o valor modal.
Esse foi também o conceito modal em duas regiões : Nordeste ( 41,2 \% ) e Sudeste ( 35,4 \% ).
O conceito 3 foi o segundo mais frequente em nível nacional ( 30,3 \%, correspondendo a 33 cursos ), sendo o conceito modal da região Sul ( 33,3 \% ), e o conceito 2, o terceiro ( 18,3 \%, correspondendo a 20 cursos ).
Houve, ainda, 12 cursos ( 11,0 \% ) que receberam o conceito 5, conceito modal na região Centro-Oeste ( 75,0 \% ), e sete cursos ( 6,4 \% ) que receberam o conceito 1.
Dos 109 cursos de Design, um ( 0,9 \% ) ficou Sem Conceito ( SC ).
A região Norte participou com quatro cursos, o correspondente a 3,7 \% do total nacional.
A distribuição dos conceitos na região Norte foi multimodal, com a mesma proporção para os conceitos 1, 2, 3 e 4 ( 25 \% ).
A região Nordeste participou com 17 cursos, o correspondente a 15,6 \% do total nacional.
Nessa região, sete cursos ( 41,2 \% em termos regionais ) obtiveram o conceito 4, o conceito modal para a região, como já comentado.
Os conceitos 1 e 2 foram atribuídos, respectivamente, a um e a três cursos ( 5,9 \% e 17,6 \%, respectivamente ).
Ainda nessa região, quatro cursos ( 23,5 \% ) foram avaliados com o conceito 3, e dois cursos ( 11,8 \% ), com o conceito 5.
Nenhum curso ficou Sem Conceito ( SC ).
Dos 48 cursos participantes da região Sudeste ( 44,0 \% do total nacional ), o conceito 4 foi o conceito modal, como já anteriormente assinalado, tendo sido obtido por 17 cursos ( 35,4 \% ).
O conceito 1 foi atribuído a dois cursos ( 4,2 \% ), e o conceito 2, a 11 cursos ( 22,9 \% ).
Receberam os conceitos 3 e 5, respectivamente, 15 cursos ( 31,3 \% ) e dois cursos ( 4,2 \% ).
Nessa região, um curso ficou Sem Conceito ( SC ).
Os 36 cursos da região Sul corresponderam a 33,0 \% do total nacional.
Os conceitos 1 e 2 foram atribuídos a três ( 8,3 \% ) e a cinco cursos ( 13,9 \% ), respectivamente.
Como já indicado, houve predominância do conceito 3, atribuído a 33,3 \% dos cursos da região ( conceito modal ), ou seja, atribuído a 12 dos 36 cursos participantes na região Sul.
Já os conceitos 4 e 5 foram atribuídos a 11 ( 30,6 \% ) e a cinco cursos ( 13,9 \% ), respectivamente.
Nenhum curso ficou Sem Conceito ( SC ) na região Sul.
Dos quatro cursos participantes na região Centro-Oeste ( 3,7 \% do total nacional ), três ( 75,0 \% em termos regionais ) receberam conceito 5, o conceito modal, como já fora destacado, e um curso ( 25,0 \% ) obteve o conceito 3.
Nessa região, nenhum curso ficou Sem Conceito ( SC ).
No Gráfico 5.1, é apresentada a distribuição cumulativa do Conceito Enade, segundo a Grande Região.
Em um gráfico de distribuição cumulativa, em cada valor de abscissa ( eixo x ), o valor apresentado no eixo das ordenadas ( eixo y ) é o valor acumulado da variável ( conceito Enade ).
Observando -se a linha vertical pontilhada, que ajuda a visualizar o percentual de conceitos acumulados até o valor 3, ou seja, conceitos 1, 2 e 3, pode -se notar que, por exemplo, a região Norte ( linha azul ) apresentou 75 \% dos cursos com conceito 3 ou abaixo desse valor.
A região Centro-Oeste ( linha laranja ), por outro lado, apresentou 25 \% dos cursos na mesma situação : com conceito 3 ou abaixo.
Destaca -se que, em algumas regiões, o conceito 5 pode não ter sido alcançado por nenhum curso.
Foi o que ocorreu na região Norte ( linha azul ), em que todos os cursos obtiveram conceitos menores ou iguais a 4, com a poligonal alcançando 100 \% no conceito 4.
Em linhas gerais, em um gráfico de distribuição cumulativa, poligonais mais à esquerda representam as Grandes Regiões com a pior distribuição, e poligonais mais à direita, as Grandes Regiões com os melhores conceitos.
A região Norte ( linha azul ) e a Sudeste ( linha vermelha ) apresentaram os cursos com os conceitos mais baixos, ao passo que a Centro-Oeste ( linha laranja ) apresentou a distribuição com os valores maiores.
Cursos em IES situadas nas regiões Sul ( linha rosa ) e Nordeste ( linha verde ) apresentaram uma situação intermediária. 5.2 CONCEITOS POR CATEGORIA ADMINISTRATIVA, POR MODALIDADE DE ENSINO E GRANDE REGIÃO Na Tabela 5.2, é apresentada a distribuição dos cursos de Design participantes do Enade/2021, por Categoria Administrativa e Modalidade de Ensino, de acordo com os conceitos por eles alcançados, segundo as Grandes Regiões brasileiras.
Quanto à Modalidade de Ensino, todos os cursos eram Presenciais.
Dos 109 cursos participantes, 26 ( 23,9 \% ) foram ministrados em Instituições Públicas, e 83 ( 76,1 \% ), em Instituições Privadas.
O conceito 4 foi o valor modal, atribuído a 36 cursos ( 33,0 \% ).
Quanto aos demais cursos, sete cursos ( 6,4 \% ) receberam o conceito 1 ; 20 cursos ( 18,3 \% ), o conceito 2 ; 33 cursos ( 30,3 \% ), o conceito 3 ; e 12 cursos ( 11,0 \% ), o conceito 5.
Apenas um curso ( 0,9 \% ) ficou Sem Conceito ( SC ). conceito 3 ; e 12 cursos ( 11,0 \% ), o conceito 5.
Apenas um curso ( 0,9 \% ) ficou Sem Conceito ( SC ).
De acordo com as informações da Tabela 5.2, em termos nacionais, entre Instituições Públicas e Privadas, observa -se que, dos 12 cursos avaliados com o conceito 5, seis eram oferecidos em IES Públicas, e seis, em IES Privadas.
Dos 26 cursos participantes de IES Públicas, o conceito 4 foi o valor modal, atribuído a 14 cursos ( 53,8 \% ).
Dos demais cursos de Instituições Públicas participantes, quatro cursos ( 15,4 \% ) foram avaliados com o conceito 3 ; um curso ( 3,8 \% ) recebeu o conceito 2 ; e um curso ( 3,8 \% ), o conceito 1.
Nessa categoria, nenhum curso ficou Sem Conceito ( SC ).
Na rede privada, o conceito modal foi o 3, atribuído a 29 cursos, o correspondente a 34,9 \% dos 83 cursos da categoria.
Dos demais cursos participantes, seis ( 7,2 \% ) receberam o conceito 1, e 19 ( 22,9 \% ), o conceito 2.
O conceito 4 foi atribuído a 22 cursos ( 26,5 \% ), e o conceito 5, a seis cursos ( 7,2 \% ).
Nessa Categoria Administrativa, um curso ( 1,2 \% ) ficou Sem Conceito ( SC ).
Na análise por região, observa -se que, na região Norte, foram quatro os cursos participantes.
As Instituições Públicas participaram com um curso ( 25,0 \% do total regional ), ao qual foi atribuído o conceito 4.
Já as Instituições Privadas participaram com três cursos ( 75,0 \% do total regional ), e desses, um foi avaliado com o conceito 1, outro com o conceito 2, e o terceiro com o conceito 3.
Nenhum curso da região Norte ficou Sem Conceito ( SC ) ou com conceito 5.
Na região Nordeste, as Instituições Públicas da região participaram com 11 cursos ( 64,7 \% do total da região ).
Desses, sete ( 63,6 \% ) obtiveram o conceito 4, o valor modal, seguidos por dois cursos ( 18,2 \% ) que receberam o conceito 3.
Um curso ( 9,1 \% ) recebeu o conceito 1, e um ( 9,1 \% ) recebeu o conceito 5.
A rede privada concentrou seis dos 17 cursos participantes, o equivalente a 35,3 \% do total da região.
Entre os cursos oferecidos por IES Privadas, no Nordeste, três cursos ( 50,0 \% ) foram avaliados com conceito 2, o conceito modal.
Dois cursos obtiveram o conceito 3, e um curso ficou com o conceito 5.
Nenhum curso ficou Sem Conceito ( SC ).
Na região Sudeste, a proporção de cursos da rede pública, 12,5 \%, correspondeu a seis dos 48 cursos participantes.
Nessa categoria e região, o conceito modal foi o 4, atribuído a cinco cursos, e um curso ficou com o conceito 3.
Dos 42 cursos oferecidos por Instituições Privadas, na região Sudeste, o conceito modal foi o 3, atribuído a 14 cursos.
Na região, 12 cursos receberam o conceito 4, e 11, o conceito 2.
Ainda nessa região, dois cursos obtiveram o conceito 1, e dois, o conceito 5.
Apenas um curso ficou Sem Conceito ( SC ).
As Instituições Públicas, na região Sul, participaram com seis cursos ( 16,7 \% dos cursos da região ), aos quais foram atribuídos os conceitos 2 ( um curso ), 3 ( um curso ) e 4 ( um curso ), além do conceito modal 5, atribuído a três.
Nenhum curso da região Sul oferecido por IES Públicas ficou Sem Conceito ( SC ) ou recebeu o conceito 1.
As Instituições Privadas concentraram 30 dos 36 cursos participantes da região Sul, 83,3 \% do total regional.
Desses, 11 cursos obtiveram conceito 3, o conceito modal.
Nessa combinação de Categoria Administrativa e Grande Região, três cursos receberam o conceito 1, quatro cursos receberam o conceito 2, outros dez, o conceito 4, e ainda dois, o conceito 5.
Nenhum curso ficou Sem Conceito ( SC ).
Na região Centro-Oeste, dois dos quatro cursos participantes eram de instituições Públicas ( 50,0 \% em termos regionais ), e ambos receberam o conceito 5.
Quanto aos dois cursos oferecidos por Instituições Privadas na região Centro Oeste ( 50,0 \% do total regional ), um curso recebeu o conceito 3, e um outro, o conceito 5.
No Gráfico 5.2, apresenta -se a distribuição do Conceito Enade, segundo a Categoria Administrativa da IES.
Os cursos em IES Públicas ( linha azul ) apresentam uma poligonal mais à direita do que os em IES Privadas ( linha verde ) e, consequentemente, uma distribuição de Conceito Enade com valores maiores. 5.3 CONCEITOS POR ORGANIZAÇÃO ACADÊMICA E GRANDE REGIÃO Na Tabela 5.3, encontra -se a distribuição dos conceitos atribuídos aos cursos participantes do Enade/2021, na Área de Design, por Organização Acadêmica, segundo as Grandes Regiões brasileiras.
Dos 109 cursos de Design participantes, 70 eram oferecidos em Universidades, 18 em Centros Universitários, 20 em Faculdades e um em CEFET/IF.
Esta distribuição corresponde a, respectivamente, 64,2 \%, 16,5 \%, 18,3 \% e 0,9 \% dos cursos.
De acordo com os dados apresentados, de todos os 12 cursos avaliados com conceito 5, nove eram vinculados a Universidades.
Os cursos vinculados a esse tipo de Organização Acadêmica tiveram o conceito 4 como conceito modal, atribuído a 25 cursos ( 35,7 \% ).
Os demais cursos vinculados a Universidades receberam os conceitos 1 ( seis cursos ), 2 ( seis cursos ), 3 ( 23 cursos ) e 5 ( nove cursos, como já mencionado ).
Um curso ficou Sem Conceito ( SC ).
Dos 18 cursos em Centros Universitários, o conceito modal foi o 3, atribuído a sete cursos ( 38,9 \% ).
Os outros cursos receberam os conceitos 2 ( cinco cursos ) e 4 ( seis cursos ).
Nenhum curso vinculado a esse tipo de Organização Acadêmica ficou Sem Conceito ( SC ) ou obteve conceito 1 ou 5.
Dos 20 cursos mantidos por Faculdades, oito ( 40,0 \% ) receberam conceito 2, o conceito modal.
O segundo conceito mais frequentemente atribuído aos cursos oferecidos por essa categoria de IES foi o conceito 4, recebido por cinco cursos.
Ainda nesse tipo de Organização Acadêmica, um curso obteve o conceito 1, três cursos, o conceito 3, e três cursos, o conceito 5.
Dos cursos oferecidos em Faculdades, nenhum ficou Sem Conceito ( SC ).
A Organização Acadêmica CEFET/IF ofereceu um curso, que obteve conceito 2.
Considerando -se, separadamente, as regiões brasileiras, verifica -se que, na região Norte, apenas as Universidades e as Faculdades tiveram cursos avaliados.
As Universidades Na região Nordeste, as Universidades participaram com 13 dos 17 cursos da Área de Design oferecidos.
Nenhum desses 13 cursos ficou Sem Conceito ( SC ).
O conceito modal 4 foi atribuído a sete cursos, e três cursos obtiveram o conceito 3.
Os três cursos restantes obtiveram os conceitos 1, 2 e 5.
Os Centros Universitários contaram com três cursos participantes na região Nordeste, dois dos quais receberam o conceito modal 2, e o terceiro curso recebeu o conceito 3.
Já as Faculdades foram representadas por um curso na região Nordeste, o qual obteve o conceito 5.
Não houve cursos avaliados nos CEFET/IF, nessa região.
Na região Sudeste, as Universidades concentraram 25 dos 48 cursos de Design da região.
Dos cursos oferecidos em Universidades, nessa região, o conceito modal foi o 4, atribuído a nove cursos.
Os demais cursos receberam os conceitos 1 ( dois cursos ), 2 ( quatro cursos ), 3 ( oito cursos ) e 5 ( um curso ), e um curso ficou Sem Conceito ( SC ).
Os Centros Universitários participaram com dez cursos na região Sudeste, dos quais cinco obtiveram o conceito modal 3.
O conceito 2 foi atribuído a dois cursos, e o conceito 4, a três cursos.
Nenhum curso ficou Sem Conceito ( SC ) nessa combinação de Organização Acadêmica e região ou recebeu os conceitos 1 ou 5.
Já as Faculdades foram representadas por 13 cursos na região Sudeste.
Nesse tipo de Organização Acadêmica, a distribuição de conceito foi bimodal : cinco cursos receberam o conceito 2, e cinco, o conceito 4.
Aos demais cursos foram atribuídos os conceitos 3 ( dois cursos ) e 5 ( um curso ).
Nenhum curso recebeu conceito 1 ou ficou Sem Conceito ( SC ).
Não houve cursos avaliados nos CEFET/IF, nessa região.
Dos 36 cursos da região Sul, 27 eram mantidos por Universidades, para os quais o conceito modal foi 3, atribuído a 11 cursos.
Os demais cursos receberam os conceitos 1 ( três cursos ), 2 ( um curso ), 4 ( oito cursos ) e 5 ( quatro cursos ).
Dos cinco cursos participantes de Centros Universitários da região Sul, nenhum recebeu os conceitos 1 ou 5.
Os conceitos restantes foram atribuídos a um curso ( conceito 2 ), um curso ( conceito 3 ) e três cursos ( conceito modal 4 ).
Quanto aos três cursos vinculados a Faculdades, na região Sul, dois receberam o conceito 2, e um recebeu o conceito 5.
O único curso oferecido por CEFET/IF nessa região obteve o conceito 2.
Nenhum curso da região Sul ficou Sem Conceito ( SC ).
Na região Centro-Oeste, todos os quatro cursos foram mantidos por Universidades.
Desses, três obtiveram o conceito 5 ( modal ), e um obteve o conceito 3.
No Gráfico 5.3, apresenta -se a distribuição do Conceito Enade dos cursos de Design, segundo a Organização Acadêmica da IES.
Os cursos em Universidades ( linha azul ) participaram com apenas um dos quatro cursos avaliados, ao qual foi atribuído o conceito 4.
As Faculdades participaram com os três cursos restantes nessa região, os quais receberam os conceitos 1, 2 e 3.
Na região Nordeste, as Universidades participaram com 13 dos 17 cursos da Área de Design oferecidos.
Nenhum desses 13 cursos ficou Sem Conceito ( SC ).
O conceito modal 4 foi atribuído a sete cursos, e três cursos obtiveram o conceito 3.
Os três cursos restantes obtiveram os conceitos 1, 2 e 5.
Os Centros Universitários contaram com três cursos participantes na região Nordeste, dois dos quais receberam o conceito modal 2, e o terceiro curso recebeu o conceito 3.
Já as Faculdades foram representadas por um curso na região Nordeste, o qual obteve o conceito 5.
Não houve cursos avaliados nos CEFET/IF, nessa região.
Na região Sudeste, as Universidades concentraram 25 dos 48 cursos de Design da região.
Dos cursos oferecidos em Universidades, nessa região, o conceito modal foi o 4, atribuído a nove cursos.
Os demais cursos receberam os conceitos 1 ( dois cursos ), 2 ( quatro cursos ), 3 ( oito cursos ) e 5 ( um curso ), e um curso ficou Sem Conceito ( SC ).
Os Centros Universitários participaram com dez cursos na região Sudeste, dos quais cinco obtiveram o conceito modal 3.
O conceito 2 foi atribuído a dois cursos, e o conceito 4, a três cursos.
Nenhum curso ficou Sem Conceito ( SC ) nessa combinação de Organização Acadêmica e região ou recebeu os conceitos 1 ou 5.
Já as Faculdades foram representadas por 13 cursos na região Sudeste.
Nesse tipo de Organização Acadêmica, a distribuição de conceito foi bimodal : cinco cursos receberam o conceito 2, e cinco, o conceito 4.
Aos demais cursos foram atribuídos os conceitos 3 ( dois cursos ) e 5 ( um curso ).
Nenhum curso recebeu conceito 1 ou ficou Sem Conceito ( SC ).
Não houve cursos avaliados nos CEFET/IF, nessa região.
Dos 36 cursos da região Sul, 27 eram mantidos por Universidades, para os quais o conceito modal foi 3, atribuído a 11 cursos.
Os demais cursos receberam os conceitos 1 ( três cursos ), 2 ( um curso ), 4 ( oito cursos ) e 5 ( quatro cursos ).
Dos cinco cursos participantes de Centros Universitários da região Sul, nenhum recebeu os conceitos 1 ou 5.
Os conceitos restantes foram atribuídos a um curso ( conceito 2 ), um curso ( conceito 3 ) e três cursos ( conceito modal 4 ).
Quanto aos três cursos vinculados a Faculdades, na região Sul, dois receberam o conceito 2, e um recebeu o conceito 5.
O único curso oferecido por CEFET/IF nessa região obteve o conceito 2.
Nenhum curso da região Sul ficou Sem Conceito ( SC ).
Na região Centro-Oeste, todos os quatro cursos foram mantidos por Universidades.
Desses, três obtiveram o conceito 5 ( modal ), e um obteve o conceito 3.
No Gráfico 5.3, apresenta -se a distribuição do Conceito Enade dos cursos de Design, segundo a Organização Acadêmica da IES.
Os cursos em Universidades ( linha azul ) apresentam uma poligonal mais à direita do que as demais, denotando uma melhor distribuição de Conceitos Enade.
No outro extremo, encontram -se os cursos mantidos por CEFET/IF ( linha laranja ), apresentando as respectivas poligonais mais à esquerda e, consequentemente, a distribuição de Conceitos Enade com os menores valores.
Os cursos mantidos por Centros Universitários ( linha verde ) e por Faculdades ( linha vermelha ) apresentam poligonais mais centrais e, consequentemente, a distribuição de Conceitos Enade com valores intermediários.

% matched lemmas: 
\end{textsample}
